\chapter{Why Now?}
\label{ch:whynow}

The mathematics of neural networks is old: the perceptron dates from 1958 and
the standard training method from the 1970s and 1980s. The question is why they
started working this well only recently. Three trends met at about the same
time, and none would have been enough alone.

\section{Three trends}

\textbf{Compute.} For decades Moore's law doubled the number of transistors on a
chip roughly every two years. Graphics processors (GPUs), built to draw
triangles, turned out to be fast at exactly the operation neural networks spend
their time on: multiplying large matrices in parallel. This made GPUs useful
for training neural networks. What changed recently is the scale of money
behind it. Training a frontier model is now out of reach for any university;
running one on hardware you can buy is not, and that gap is what
\Cref{ch:localfirst} is about. For running models, the limit is memory rather
than processor speed (\Cref{ch:hardware}).

\textbf{Algorithms.} \textcite{vaswani2017attention} introduced the
\term{transformer} in 2017, using \term{attention} to combine information from
different positions in a sequence. During training, many positions can be
processed in parallel, making better use of GPUs than processing them one
after another. When predicting the next token, the model is prevented from
looking ahead at the answer. Generating new text still proceeds token by
token in this design (\Cref{fig:parallel}). The original experiments concerned translation and
sentence structure; they did not establish ecological expertise
\parencite{vaswani2017attention}. The design continues to develop:
\textcite{kimi2026attnres}, for instance, change how information from earlier
layers is combined.

\begin{figure}[htbp]
\centering
\fig{ch02-parallel-sequential}
\caption{Training and generation in a transformer. During training every
position predicts its next token at once, without seeing the tokens after it;
when generating, the model re-reads the whole sequence and adds one token at a
time.}
\label{fig:parallel}
\end{figure}

\textbf{Data.} Pretraining consumed the accessible text of the internet. For
ecology the equivalent resource is domain data --- curated, standardised, with
known provenance --- and it is the one of the three we own. It decides whether
any of this produces science rather than plausible text, and it is the subject
of Part~II.

\section{A race, and who is running it}

The framing has moved from products to national industrial strategy
\parencite{openai2026industrialpolicy}. The AI Index \parencite{aiindex2026}
gives the scale. Private investment in AI worldwide came to \$344.7 billion in
2025, more than double the year before, and the United States attracted about
23 times more than the next country. The US produced 59 notable models in 2025,
and China produced 35. China leads in publications and patents. The top Chinese model
trailed the top American one by under 3 per cent on a public leaderboard in
March 2026. Europe produces few frontier models, yet across 25 countries more
people trust the EU than the US or China to regulate AI.

Open and closed models are racing too. The best closed model led the best
open-weight one by 15 per cent in 2023, by under 1 per cent in 2024 and by 3.4
per cent in March 2026 \parencite{aiindex2026} (\Cref{fig:gap}). Open weights are close enough
to the frontier that running them yourself is a real option, which is the
subject of Part~IV.

\begin{figure}[htbp]
\centering
\fig{ch02-open-closed-gap}
\caption{The lead of the best closed model over the best open-weight one on a
public leaderboard, in per cent, from 2023 to March 2026
\parencite{aiindex2026}.}
\label{fig:gap}
\end{figure}

The money is also reaching biology directly. In a commentary,
\textcite{palermo2026acquirers} examines Anthropic's acquisition of Coefficient
Bio as evidence of a potential new class of biotech buyer. Her distinction is
between buying a drug candidate and acquiring models, biological datasets and
the team able to develop them. She expects further acquisitions and
partnerships; this is a forecast, rather than an established industry pattern.
For biodiversity infrastructures, the question is what access to valuable
data will mean for the people and institutions that curate them
(\Cref{ch:aireadydata}).

\section{Is this a bubble?}

Two questions are being asked at once: whether the models work, and whether the
money comes back. The capability is real and measured. The AI Index reports
that frontier companies' revenue is growing at historically fast rates,
but spending on compute is rising as fast. One cloud provider alone reported
over \$150 billion of capital spending in 2025
\parencite{aiindex2026}. Whether that returns a profit is a question about
business models, not about whether the models work.

Two effects reach research budgets either way. Demand for memory from data
centres has, in the author's experience of buying hardware, pushed up the price
of RAM and storage for everyone. And infrastructure built on the assumption of
cheap API access becomes fragile if prices change. Keeping some capability on
your own hardware is partly insurance against both.

\subsection{Hype}

Ewan Birney, director of EMBL's European Bioinformatics Institute, compares the
moment to the genomics boom he worked in: real scientific change, bold companies
and wild predictions \parencite{birney2026genomics}. He calls the prediction ``we
will cure all human diseases in five years'' ``hubris'': many diseases are not even well defined, and proving that a cure
works is slow, laborious
work. His advice is to drop the fantastical claims and make discoveries. He
points out that genomics dealt with its ethical problems (privacy, gene patents,
the risk of eugenic thinking) through research, regulation and law. The same is
possible here.

Conservation has its own version of this. \textcite{sandbrook2024hype} argues
that conservation AI has met too much techno-optimism and too little attention
to its risks. These include use by bad actors, biased training data, harm to the
sector's staff and skills, and a high demand for energy, land and water. Estimates
of that demand per query differ widely. One study put an older model at half a
litre of water for every 10 to 50 responses \parencite{li2023thirsty}. Google
measured 0.26 millilitres for the median prompt to its own assistant
\parencite{elsworth2025google}.

\section{Data, and who paid for it}

When the open web ran out, the search for text moved to books. Court documents
in a US court case revealed an Anthropic memo describing ``Project Panama''
as ``our effort to destructively scan all the books in the world''
\parencite{james2026anthropicbooks}. Books were bought in bulk, their spines cut
off, the pages scanned and the rest thrown away. Destructive scanning is faster and cheaper
than careful scanning, and buying a book avoids negotiating a licence.
Volunteers have started racing to scan rare books non-destructively before they
enter that pipeline \parencite{annasarchive2026scanning}. (This book was written
with tools from the same company; the point stands.)

For biodiversity this is not abstract. The Biodiversity Heritage Library holds
taxonomic literature back to the eighteenth century --- original species
descriptions, type designations, expedition reports --- much of it in one or a
few physical copies. It is both the foundation of nomenclature and an attractive
training corpus, and the institutions holding it are not funded to defend it.

\begin{keypoint}
Ecology's data commons was built by people, over decades, to be shared.
``AI-ready'' describes what that work made possible. It does not say who did it,
or on what terms it is now being consumed.
\end{keypoint}
